\section{The required chessboard homology}\label{sec:topology}

This appendix proves the full homological input, including its chain-level
prerequisites. All coefficients are in \(\F\). The identifiers H01--H13
match the dependency route of the formalization; they are included only to
make correspondence with the linked sources easy. The chessboard bound is
the homological consequence of the stronger connectivity theorem of
Bj\"orner, Lov\'asz, Vre\'cica, and \v{Z}ivaljevi\'c
\cite[Theorem 1.1]{BLVZ94}. We follow its star-cover strategy, with an
explicit finite chain proof of the cover lemma. We do not assume that
theorem or a spectral-sequence convergence result as a black box.

\subsection{Augmented chains and their elementary operations}

A finite augmented simplicial complex \(K\) is a downward-closed family
of subsets of a finite vertex set \(V\), containing the empty face.
It need not have a vertex. Write \(C_k(K)\) for coefficient functions
on its faces with exactly \(k\) vertices, extended by zero to all subsets
of \(V\). Thus the cell count \(k\) is reduced homological degree
\(k-1\): \(C_0(K)=\F\) is the augmentation group.
Write \(\operatorname{Exact}(K,k)\) when every \(k\)-cell cycle has a
\((k+1)\)-cell filler, and \(\operatorname{Acyclic}(K,n)\) when this
holds for every \(0\le k<n\). For \(n\ge1\), this says reduced homology
vanishes from degree \(-1\) through \(n-2\); for \(n=0\) it is vacuous.
In particular, \(\operatorname{Exact}(K,0)\) means that \(K\) has a vertex.

\begin{proposition}[H01: supported augmented boundary]\label{prop:h01}
\lean{claims/AugmentedChains.lean}
For \(k\ge1\), the formula
\[
 (\partial c)(T)=\sum_{v\notin T}c(T\cup\{v\})
\]
defines a linear map \(C_k(K)\to C_{k-1}(K)\). Its restriction from
\(C_1(K)\) to \(C_0(K)\) is the sum of vertex coefficients.
The outgoing boundary of \(C_0(K)\) is zero.
The supported-coefficient model is linearly equivalent to functions on
exact-size faces, including for chessboard complexes.
\end{proposition}
\begin{proof}
A nonzero summand requires \(T\cup\{v\}\in K\) and cardinality \(k\).
Downward closure gives \(T\in K\), and \(v\notin T\) gives
\(|T|=k-1\). Finite sums are linear. At count zero every inserted face
is nonempty, so its coefficient vanishes. Restrict a supported function
to exact-size faces and invert by extension by zero. These maps are
linear and mutually inverse.
\end{proof}

\begin{lemma}[H02: the boundary squares to zero]\label{lem:h02}
\lean{third-party-claims/AugmentedBoundarySquared.lean}
The augmented boundary satisfies \(\partial^2=0\) at every cell count.
\end{lemma}
\begin{proof}
At a face \(T\), the double sum inserts two distinct vertices outside
\(T\). Pair the ordered insertion \((u,v)\) with \((v,u)\).
They contribute the same coefficient, twice, hence zero in \(\F\).
The outgoing augmentation boundary is already zero, so this includes
all low-count cases.
\end{proof}

\begin{lemma}[H03: injective chain maps and relabeling]\label{lem:h03}
\lean{claims/AugmentedChainMaps.lean}
\lean{claims/AugmentedSubcomplexRelabeling.lean}
An injective vertex map carrying faces of \(K\) to faces of \(L\)
induces a degree-preserving linear pushforward commuting with boundary.
If it gives a bijection on the faces of the two complexes, this is a
chain isomorphism and transports cycle fillings in both directions.
In particular this applies to subcomplex inclusions, chessboard transpose,
and relabeling a residual board. It also identifies the supported boundary
with the boundary on exact-size chessboard faces.
\end{lemma}
\begin{proof}
On an image face \(f(S)\), give the pushforward coefficient \(c(S)\),
and assign zero off the range. Injectivity preserves cardinality and
makes the definition unambiguous. On a face outside the range all boundary
terms vanish. On an image face, insertions in the image correspond
bijectively to insertions in the source, and all other terms vanish.
Thus pushforward commutes with boundary. A face bijection supplies the
inverse coefficient map. Swapping row and column coordinates preserves
precisely the matching condition; residual relabeling has the same property.
Restriction to exact-size faces and extension by zero identify the two
boundary formulas term by term, including augmentation.
\end{proof}

\begin{lemma}[H04: explicit cone contraction]\label{lem:h04}
\lean{claims/AugmentedCone.lean}
Suppose \(v\) is a cone vertex of \(K\), meaning that
\(S\in K\) implies \(S\cup\{v\}\in K\). Define
\[
 (s_vc)(U)=\begin{cases}c(U\setminus\{v\})&v\in U,\\0&v\notin U.\end{cases}
\]
Then \(s_v:C_k(K)\to C_{k+1}(K)\) and
\(\partial s_v+s_v\partial=\mathrm{id}\). Consequently
\(\operatorname{Exact}(K,k)\) holds for every \(k\ge0\).
\end{lemma}
\begin{proof}
The support assertion follows by adjoining the cone vertex.
If \(v\notin T\), the only nonzero term in \((\partial s_vc)(T)\)
is insertion of \(v\), giving \(c(T)\), and \((s_v\partial c)(T)=0\).
If \(v\in T\), all insertions outside \(T\) occur once in each of
\(\partial s_vc\) and \(s_v\partial c\) and cancel; the remaining
insertion of \(v\) into \(T\setminus\{v\}\) gives \(c(T)\).
For a cycle, the identity gives \(\partial(s_vc)=c\).
\end{proof}

\begin{lemma}[H05: the boundary of a simplex]\label{lem:h05}
\lean{third-party-claims/SimplexBoundary.lean}
The simplex on a nonempty finite vertex set is exact at every cell count.
If that set has size \(b\ge1\), its boundary complex, consisting of all
proper subsets, is exact at every \(k\) with \(k+1<b\).
Equivalently it is \(\operatorname{Acyclic}(\partial\Delta,n)\) whenever
\(n<b\). No top-dimensional vanishing is asserted.
\end{lemma}
\begin{proof}
The full simplex is a cone at any vertex. Contract a \(k\)-cycle in its
boundary inside the full simplex using Lemma~\ref{lem:h04}. Its filler
has \(k+1<b\) vertices on every supported face, so all those faces are
proper subsets and the filler already lies in the boundary complex.
\end{proof}

\subsection{The finite homological cover argument}

Let \((L_i)_{i\in I}\) be finitely many subcomplexes of \(K\). Put
\(L_J=\bigcap_{i\in J}L_i\), with \(L_\varnothing=K\).
Its vertex-nonempty nerve is the complex
\[
 \mathcal N=\{\varnothing\}\cup
   \{J\subseteq I:L_J\text{ has at least one vertex}\}.
\]
An empty-face-only intersection does not give a nonempty nerve face.
The family covers \(K\) if every face belongs to some \(L_i\).

\begin{lemma}[H06: cover double chains and horizontal exactness]\label{lem:h06}
\lean{claims/FiniteComplexCover.lean}\lean{claims/CoverDoubleComplex.lean}
Let \(D_{a,b}\) consist of double coefficients \(d(J,S)\) supported where
\(|J|=a\), \(|S|=b\), \(S\in L_J\), and \(J\in\mathcal N\).
Insertion boundaries \(H\) in \(J\) and \(V\) in \(S\) lower the
respective cell counts, with zero outgoing maps at count zero, and satisfy
\[
 H^2=V^2=0,\qquad HV=VH,\qquad(H+V)^2=0.
\]
If the \(L_i\) cover \(K\), then for \(b>0\) every horizontally closed
\(d\in D_{a,b}\) has a horizontal filler in \(D_{a+1,b}\).
\end{lemma}
\begin{proof}
Removing an index preserves intersection membership and nerve membership.
Removing a vertex preserves membership in every complex, while the nerve
condition on \(J\) remains unchanged. These facts prove the support claims.
The square-zero identities are Lemma~\ref{lem:h02}; both \(HV\) and
\(VH\) sum the same coefficients \(d(J\cup\{i\},S\cup\{v\})\), so
interchanging finite sums proves commutation. Characteristic two gives
the total square-zero identity.

For each nonempty supported face \(S\), choose an index \(i(S)\) whose
member contains \(S\), and contract the index coefficient
\(J\mapsto d(J,S)\) at \(i(S)\), using the ambient formula of
Lemma~\ref{lem:h04}. Call the result \(u(J,S)\), and use zero on other
\(S\). Its nonzero coefficients have \(S\) in every indicated cover
member, including the newly chosen one. Because \(S\) contains a vertex,
the enlarged index set really belongs to the nerve. Thus
\(u\in D_{a+1,b}\), and the contraction identity gives \(Hu=d\).
The condition \(b>0\) is essential: the \(b=0\) edge is the nerve
chain complex and is not automatically exact.
\end{proof}

\begin{lemma}[H07: homological cover lemma]\label{lem:nerve}
\lean{third-party-claims/HomologicalCover.lean}
Suppose \((L_i)\) covers \(K\). Let \(n\ge0\), and assume
\(\operatorname{Acyclic}(\mathcal N,n)\). Suppose also that for every
nonempty \(J\in\mathcal N\) and every \(r\ge0\) with \(|J|+r\le n\),
\(\operatorname{Exact}(L_J,r)\) holds. Then
\(\operatorname{Acyclic}(K,n)\).
In reduced degrees, putting \(n=q+2\), this is the usual assertion that
a \(q\)-acyclic nerve and \((q-|J|+1)\)-acyclic intersections give a
\(q\)-acyclic union, including the augmentation case \(q=-1\).
\end{lemma}
\begin{proof}
We use the finite double coefficients of Lemma~\ref{lem:h06}.
First note two filling facts. On the \(b=0\) edge, \(d(J,\varnothing)\)
is a nerve chain; nerve exactness fills horizontally at counts \(a<n\),
and the filler is vertically closed. For \(a>0\), a vertically closed
\(d\in D_{a,b}\) with \(a+b\le n\) fills vertically in \(D_{a,b+1}\):
for each \(J\) of size \(a\) use the assumed filling in \(L_J\), and
use zero for the other index sets. Nonempty \(J\) ensures that all filler
faces lie in \(K\) as well as the required intersection.

We claim that whenever \(a+b<n\) and \(Hd=Vd=0\), there is
\(u\in D_{a+1,b}\) with \(Hu=d\) and \(Vu=0\).
Induct on \(b\). The case \(b=0\) is the nerve-edge fact.
For \(b>0\), horizontal exactness first gives \(x\in D_{a+1,b}\)
with \(Hx=d\). Then \(Vx\in D_{a+1,b-1}\) is closed in both directions.
Induction gives \(y\in D_{a+2,b-1}\) with \(Hy=Vx\), \(Vy=0\).
The vertical filling fact gives \(z\in D_{a+2,b}\) with \(Vz=y\),
since \((a+2)+(b-1)=a+b+1\le n\). Set \(u=x+Hz\). Then
\[
 Hu=Hx+H^2z=d,\qquad Vu=Vx+HVz=Vx+Hy=0.
\]
This proves the claim by finite induction.

Finally embed a \(k\)-cell cycle \(c\) of \(K\), \(k<n\), as
\(d(\varnothing,S)=c(S)\) in \(D_{0,k}\), with other coefficients zero.
It is closed in both directions. The claim gives \(u\in D_{1,k}\)
with \(Hu=d\), \(Vu=0\). The singleton-intersection filling gives
\(z\in D_{1,k+1}\) with \(Vz=u\), because \(1+k\le n\).
The chain \(f(S)=(Hz)(\varnothing,S)\) belongs to \(C_{k+1}(K)\), and
\(\partial f=(VHz)(\varnothing,\cdot)=(HVz)(\varnothing,\cdot)=c\).
Thus every required cycle fills, including \(k=0\).
\end{proof}

\subsection{Chessboard stars, arithmetic, and assembly}

The chessboard complex \(\Delta_{a,b}\) has vertex set \([a]\times[b]\)
and faces the matchings, namely sets with no repeated row or column.
The closed star of a vertex \(v\) consists of faces \(S\) such that
\(S\cup\{v\}\) is a face. It is a subcomplex and a cone at \(v\).

\begin{lemma}[H08: first-row star cover]\label{lem:h08}
\lean{third-party-claims/ChessboardStarCover.lean}
If \(a+1\le b\), the stars of \((1,j)\), \(j\in[b]\), cover
\(\Delta_{a+1,b}\).
\end{lemma}
\begin{proof}
A matching using row one belongs to the star of its cell in that row.
Otherwise it uses at most \(a<b\) columns, leaving some column \(j\)
free. It can then be extended by \((1,j)\).
\end{proof}

\begin{lemma}[H09: intersections and residual-board chain maps]\label{lem:h09}
\lean{third-party-claims/ChessboardStarIntersections.lean}
The intersection of \(t\ge2\) distinct first-row stars in
\(\Delta_{a+1,b}\) consists exactly of matchings avoiding row one and
those \(t\) columns. It is chain-isomorphic to \(\Delta_{a,b-t}\),
including when \(t=b\) and the latter has only its empty face.
\end{lemma}
\begin{proof}
A matching in two different stars cannot use row one, since its row-one
cell would have to be both designated cells. It cannot use a designated
column in another row either, since extension by that star's cell would
fail. Conversely, avoiding all those rows and columns permits every
required extension. Relabel the remaining rows and columns and use
Lemma~\ref{lem:h03} to obtain the chain isomorphism.
\end{proof}

\begin{lemma}[H10: the star nerve and singleton fillings]\label{lem:h10}
\lean{third-party-claims/ChessboardStarNerve.lean}
For \(a\ge1\), \(b\ge2\), the nerve of the first-row stars of
\(\Delta_{a+1,b}\) is the boundary of the simplex on its \(b\) columns.
It is \(\operatorname{Acyclic}(\mathcal N,n)\) for every \(n<b\).
Every individual star is exact at every cell count.
\end{lemma}
\begin{proof}
A proper set of columns leaves an unused column and another row; their
cell belongs to every indicated star. The full set of stars has no
common vertex, by Lemma~\ref{lem:h09}. Thus the nerve is exactly the
proper-subset complex. Apply Lemma~\ref{lem:h05}; each singleton star
is a cone, so Lemma~\ref{lem:h04} fills its cycles.
\end{proof}

\begin{lemma}[H11: induction parameters]\label{lem:h11}
\lean{claims/ChessboardParameterArithmetic.lean}
For natural \(a,b\), put
\[
 \nu(a,b)=\min\left\{a,b,\left\lfloor\frac{a+b+1}{3}\right\rfloor\right\}.
\]
It is symmetric, bounded by each side, zero if a side is zero, and at least
one when both sides are positive. For \(b\ge1\), \(\nu(1,b)=1\).
If \(2\le a\le b\), then \(\nu(a,b)\le b-1\), and for \(2\le t<b\),
\[
 \nu(a,b)+1\le\nu(a-1,b-t)+t,
 \qquad \min\{a-1,b-t\}<\min\{a,b\}.
\]
For \(s\ge1\), \(N\ge2s-1\) gives \(\nu(s,N)=s\).
Cell counts \(0\le k<n\) correspond exactly to reduced degrees
\(-1\le k-1\le n-2\), with subtraction interpreted in the integers.
\end{lemma}
\begin{proof}
Symmetry and the basic bounds follow directly from the minimum.
If \(a<b\), \(\nu\le a\le b-1\); if \(a=b\ge2\), then
\(\lfloor(2b+1)/3\rfloor\le b-1\).
For the recursive bound write \(\nu=\nu(a,b)\). Each smaller-board
entry is at least \(\nu-t+1\):
\[
 a-1\ge\nu-t+1,\qquad b-t\ge\nu-t+1,
 \qquad a+b-t\ge3\nu-1-t\ge3(\nu-t+1).
\]
The last inequality uses \(t\ge2\). Taking the smaller minimum proves
the bound; if the right side is negative the comparison is automatic.
The induction measure decreases because \(a-1<a=\min\{a,b\}\).
Finally \(N\ge2s-1\) implies \(N\ge s\) and
\(\lfloor(s+N+1)/3\rfloor\ge s\), proving the narrow-board identity.
The degree translation is the integer inequality obtained by subtracting one.
\end{proof}

\begin{theorem}[H12: homological chessboard bound {\cite{BLVZ94}}]
\label{thm:chessboard}\lean{third-party-claims/ChessboardHomology.lean}
For \(a,b\ge1\), every \(k\)-cell cycle of \(\Delta_{a,b}\) has a
\((k+1)\)-cell filler whenever \(0\le k<\nu(a,b)\). Equivalently,
\(\widetilde H_j(\Delta_{a,b};\F)=0\) for
\(-1\le j\le\nu(a,b)-2\).
\end{theorem}
\begin{proof}
Induct strongly on \(\min\{a,b\}\), transposing by Lemma~\ref{lem:h03}
so that \(a\le b\). For \(a=1\), \(\nu=1\), and a vertex fills the
augmentation generator. For \(a\ge2\), cover by the first-row stars.
Their nerve is acyclic through count \(\nu\), since \(\nu\le b-1\).
Consider a nonempty nerve face \(J\) of size \(t\) and a count \(r\)
with \(t+r\le\nu\). For \(t=1\) the star is a cone. Otherwise
\(2\le t<b\), and its intersection is \(\Delta_{a-1,b-t}\), with both
sides positive. Write \(\mu=\nu(a-1,b-t)\).
Lemma~\ref{lem:h11} gives \(\nu+1\le\mu+t\), hence \(r<\mu\).
The smaller induction measure and the chain isomorphism of
Lemma~\ref{lem:h09} supply the required intersection filling.
Lemma~\ref{lem:nerve} now fills every \(k<\nu\) cycle of the original board.
\end{proof}

\begin{corollary}[H13: the matching-moment filling interface]
\label{cor:chessboard-filling}\lean{third-party-claims/ChessboardFillingProof.lean}
Let \(s\ge2\), \(N\ge2s-1\). Every cycle on exact-size
\((s-1)\)-faces of \(\Delta_{s,N}\) is the boundary of a chain on
exact-size \(s\)-faces. In reduced homology,
\[
 \widetilde H_{s-2}(\Delta_{s,N};\F)=0.
\]
For \(s=2\), the cycle hypothesis includes zero augmentation.
\end{corollary}
\begin{proof}
Lemma~\ref{lem:h11} gives \(\nu(s,N)=s\).
Apply Theorem~\ref{thm:chessboard} at cell count \(s-1<s\), then use the
exact-face boundary equivalence of Lemma~\ref{lem:h03}. This is precisely
the boundary equation used in Theorem~\ref{thm:moments}.
\end{proof}
